\documentclass[10pt,a4paper]{article}

\usepackage[utf8]{inputenc}
\usepackage[cyr]{aeguill}
\usepackage[francais]{babel}
\usepackage{hyperref}
% Les packages
%\usepackage[french]{babel}
\usepackage{fancybox}
\usepackage{graphics}
\usepackage{color}
\usepackage{latexsym}
\usepackage{amsmath,amsthm,amssymb}
\usepackage[plain]{fullpage} 
\usepackage{verbatim}
\usepackage{times,epsfig}
\usepackage{listings}
\usepackage{multirow}
\usepackage{amsmath,amsthm,amssymb,amscd,txfonts,bbm} % RAJOUT ICI !!!
\usepackage{graphics,epsfig} % RAJOUT ICI !!!
%\usepackage{algorithmic}
\usepackage[final]{pdfpages} 
\usepackage{fancyvrb}
\usepackage{enumerate}
\usepackage{listingsutf8}
\usepackage{comment}
\usepackage{minted}

% small et verbatim: to be coherent with \file
\newenvironment{smallverbatim}%
{\endgraf\small\verbatim}{\endverbatim}

\newenvironment{qv}
{\quote\Verbatim}
{\endquote\endVerbatim}

\usepackage{lastpage}

\newcommand{\terme}[1]{\texttt{#1}}
\newcommand{\termSet}[1]{\{\texttt{#1}\}}
\newcommand{\guill}[1]{<<\,#1\,>>}
% Macro pour l'inclusion de figure  
\newcommand{\fig}[2]{%
  \begin{figure*}[hbt]
   \begin{center}
      \includegraphics[width=6cm]{#1}
  \caption{\label{#1}  #2}
  \end{center}
 \end{figure*}
}

\newcommand{\figTex}[2]{%
  \begin{figure*}[hbt]
 \centerline{\input{#1.pstex_t}}
  \caption{\label{#1}  #2}
 \end{figure*}
}


% Macro commentaires
\definecolor{turquoise}{rgb}{0.20 ,0.60, 0.60}
\definecolor{vert}{rgb}{0.40,  0.60 ,0.00}
\definecolor{violet}{rgb}{0.80 , 0.00 , 0.80}
\definecolor{bleu}{rgb}{0.20,  0.20,  0.80}
\definecolor{rose}{rgb}{1.00,  0.20,  0.40}

\lstset{basicstyle=\footnotesize,showstringspaces=false,language=XML,breaklines=false,columns=fullflexible,frame=shadowbox, captionpos=b}

\newenvironment{solution}[0]{
~\\ \color{red} \textbf{Solution : }\color{black}}{%
}
%\excludecomment{solution}


\newcommand{\cle}[1]{\textbf{#1}}
\def\fk#1{\textit{#1}}
\def\tble#1{\textit{#1}}

\parindent=0pt           
                             
\usepackage{url}

\newtheorem{Exercice}{Definition}


\begin{document}

\title{Sujet d'examen UE NFE204 - Session 1}
\author{Année universitaire 2024-2025}
\date{Examen première session:  28 janvier 2025}

\maketitle

\begin{center}
{\bf \Large Consignes}\\
\begin{tabular}{| p{15cm} |}
\hline
Dur\'ee de l'examen : 3h00; Calculatrice autoris\'ee\,; Documents autorisés: 5 pages recto-verso de note personnelles.
Ce sujet contient \pageref{LastPage} pages.\\
\hline
\end{tabular}
\end{center}


Nous gérons un site de commerce électronique qui recense des restaurants
et permet à des internautes de les noter. Au départ,  nous utilisons
un système relationnel, dont voici le schéma (très simplifié).

\begin{itemize}
  \item Restaurant   (\textbf{id}, nom, catégorie)
  \item Internaute (\textbf{id}, nom)
  \item Notation (\textbf{\textit{idRestau}}, \textbf{\textit{idInternaute}}, note, commentaire)
\end{itemize}

\vspace*{0.2cm}

La catégorie du restaurant indique son niveau de prix, entre * et ****. 
Voici  un tout petit échantillon de la base.

\vspace*{0.2cm}
\begin{small}
\begin{minipage}[b]{4cm}
\begin{tabular}{l|l|l}
\cle{id} & nom & catégorie\\
\hline
20 & Bobo & ** \\
21 & Le régal & ****\\
\hline
\end{tabular}

\centerline{La table \tble{Restaurant}}
\end{minipage} \hspace*{0.5cm} \
\begin{minipage}[b]{2cm}
\begin{tabular}{l|l}
\cle{id} &  nom  \\
\hline
100 & Marie \\
101 & Adrien \\
\hline
\end{tabular}
\centerline{La table \tble{Internaute}}
\end{minipage} \hspace*{1cm} \
\begin{minipage}[b]{3cm}
\begin{tabular}{l|c|l|l}
\cle{\fk{idRestau}} & \cle{\fk{idInternaute}}  & note & commentaire \\
\hline
20  & 100  & 1 & Nul\\
20 & 101  & 5 & Parfait\\
21 & 101  & 3 & Moyen\\
\hline
\end{tabular}

\centerline{La table \tble{Notation}}
\end{minipage}
\end{small}

On veut implanter  un algorithme de filtrage collaboratif. Un des calculs consiste notamment
à déterminer la note moyenne, en fonction de
certains critères comme la catégorie du restaurant ou l'internaute.
 
\section{Documents structurés (4 pts)}

\begin{enumerate}
  \item Donnez une représentation de l'échantillon ci-dessus sous la forme d'une collection
      de documents JSON, centrée sur les entités ``Restaurant''. 
      Nous l'appellerons représentation A.
  \item Proposez une autre représentation, centrée sur les entités ``Internaute''.  
      Nous l'appellerons représentation B.
\end{enumerate}

\begin{solution}
Vu de nombreuses fois en cours. Voici pour la représentation A.
\begin{minted}{json}
   [{
      "_id": 10,
      "nom": "Bobo",
      "catégorie": "**",
      "notes": [{"id_inter": 100, "nom": "Marie", 
                     "commentaire": "Nul",  "note": 1},
                   {"id_inter": 101, "nom": "Adrien", 
                     "commentaire": "Parfait",  "note": 5}
                ]
     },
    {
    "_id": 21,
    "nom": "Le régal",
    "catégorie": "****",
    "notes": [{"id_inter": 101, "nom": "Adrien", 
                     "commentaire": "Moyen",  "note": 3}
                ]
  ]
\end{minted}
\end{solution}

\section{Modélisation Cassandra (5 pts)}

On veut proposer un schéma Cassandra pour explorer cette base. 
Rappelons d'abord ce qu'est une clé composite en Cassandra\,: 
\textit{A compound primary key consists of the partition key and the clustering key. The partition key determines 
which node stores the data. Rows 
for a partition key are stored in order based on the clustering key.} 

On crée les tables suivantes, correspondant à une modélisation Cassandra de 
nos documents de type A\,:
\begin{minted}{sql}
create table Restaurant (catégorie text, id_restau text, nom restau text, adresse text,
                     primary key (catégorie, id_restau));

create type Internaute (id_internaute integer, nom text, prénom text);                 

create table Notation (id_restau text,
                       internaute frozen<Internaute>,
                       note integer,
                       commentaire text,
                       primary key (id_restau, internaute));
\end{minted}

\begin{enumerate}
    \item Commentez chacune des affirmations suivantes, indiquez si elles
sont vraies ou fausses et pourquoi.

\begin{enumerate}
  \item Les lignes de \texttt{Restaurant} et les lignes de \texttt{Notation} 
     sont sur le même serveur
   \item Les restaurants d'une même catégorie sont sur le même serveur
   \item Le choix de \texttt{catégorie} comme clé de partitionnement est judicieux.
   \item Les notations d'un restaurant sont sur un seul serveur mais distribuées
        de manière aléatoire dans le stockage local.
   \item Les notations d'un même internaute sont sur un seul serveur et triées sur l'identifiant du restaurant.  
\end{enumerate}  

  \item Expliquez maintenant quelles sont les requêtes et 
     opérations nécessaires pour calculer la moyenne des notes 
     d'une catégorie de restaurant (par exemple les '***'). 
      Discutez de l'efficacité.     
  \item Proposez un second schéma avec une table \texttt{Internaute}
       correspondant à nos documents de type B
      (première question du sujet).
\end{enumerate}

\begin{solution}
\begin{enumerate}
\item
\begin{enumerate}
  \item Non puisqu'elles n'ont pas la même clé de partitionnement.
   \item Oui, la catégorie est la clé de partitionnement
   \item Non car il n'y a que 5 catégories, donc un partitionnement limité à 5 fragments.
   \item Oui sur un seul serveur mais de plus le stockage est contigu.
   \item Pas du tout car l'internaute n'est pas clé de partitionnement. 
\end{enumerate} 
  \item Une première requête va récupérer tous les restaurants 
   '***'. C'est un parcours séquentiel sur un seul serveur.
\begin{minted}{sql}
    select * from Restaurant where catégorie='***';
\end{minted}

    Ensuite, pour chaque identifiant de restaurant on
    récupère les notes. Là aussi, parcours séquentiel sur un seul serveur
    (mais autant de requêtes que de restaurants).
\begin{minted}{sql}
    select note from Notation where id_restau=<id>;
\end{minted}
     On a toutes les notes, on peut calculer la moyenne au fur et à mesure.
   \item On peut tout mettre dans une seule table,
  en intégrant les notes sous forme de liste (en espérant qu'il
  n'y en aura pas trop).

\begin{minted}{sql}
create type Note (id_restau integer, note integer, commentaire text);                 

create table Internaute (id_internaute text,
                         nom text,
                         notes list(Note),
                         primary key (id_internaute));
\end{minted}

\end{enumerate}
\end{solution}

\section{Spécification de calculs distribués (4 pts)}

On suppose que les informations sont disponibles sous la
forme d'une collection de documents JSON structurés 
selon la représentation A (cf. première question du sujet).
     
\begin{enumerate}
  \item Indiquez
    le traitement MapReduce qui, appliqué à cette collection de documents,
    permet de produire les lignes à insérer dans la 
    table Cassandra \texttt{Notation}
    donné dans la section précédente. 

\item  Indiquez
    le traitement MapReduce qui, appliqué à cette collection de documents,
    permet de produire les lignes à insérer dans la 
    table Cassandra \texttt{Internaute} dont vous avez
    proposé un schéma dans la section précédente. Vous pouvez aussi 
    expliquer comment ce traitement permet de produire la  
    représentation B à partir de la représentation A.
   
   \item La phase de Reduce (et donc celle de \textit{shuffle}) 
   est-elle  nécessaire dans les deux cas\,? Expliquer.
\end{enumerate}

Pour spécifier les calculs MapReduce, donner les fonctions de Map et de Reduce, en javascript,
en pseudo-code, au pire en langage naturel en étant le plus précis possible.

\begin{solution}

\begin{enumerate}
  \item Pour la fonction de Map, on reçoit un document représentant un
     restaurant. Il faut boucler  sur les notes, et émettre 
     une paire clé-valeur avec comme clé l'identifiant du restaurant
     et l'internaute. On peut en fait directement insérer
     dans la table \texttt{Notation} dès le Map puisqu'aucun 
     regroupement n'est nécessaire.
  \item Pour le second en revanche, le Reduce est indispensable. Le Map 
  émet, dans une boucle sur les notes, des paires dont la clé est 
  l'identifiant de l'internaute. Le \textit{shuffle} regroupe 
  les paires de même clé ce qui permet de constituer la liste des notes, 
  et le Reduce peut alors insérer 
  une ligne pour chaque groupe dans la table \texttt{Internaute}. 
\end{enumerate}

\end{solution}
\section{Recherche d'information (4 pts)}

On ajoute dans les  documents des commentaires sur les restaurants, désignés
par $c_1, ..., c_5$. On veut maintenant
créer un index plein texte pour rechercher des commentaires. On va travailler sur l'échantillon
suivant.

\vspace*{0.3cm}

\begin{tabular}{lp{15cm}}
     $c_1$& Plats médiocres, bonne carte des vins, accueil cordial.\\
     $c_2$& Le vin est moyen, les desserts sont mauvais, mais le décor est agréable.\\
     $c_3$& En général, le vin est trop cher au restaurant, mais ici la carte
         des vins est complète et de qualité.\\
     $c_4$& Passons sur le décor d'un goût douteux. Les plats et desserts sont délicieux,
          c'est tout ce qui compte après tout.\\
     $c_5$& Plat, dessert et un verre de vin pour un prix correct. Le plat du jour
             est toujours intéressant.\\ 
\end{tabular}

\bigskip

On va se limiter au vocabulaire (\texttt{vin, plat, décor, dessert}).
(NB: on suppose
     une phase préalable de normalisation qui élimine les pluriels, majuscules, etc.)

\begin{enumerate}
   \item Donnez une matrice d'incidence contenant les tf, 
      et un tableau donnant les idf (sans le $\log$), pour les  termes précédents. Placez
      les termes en ligne, les commentaires en colonne.

\begin{solution}
\begin{center}
\begin{tabular}{l|c|c|c|c|c}
            & \textbf{c1} & \textbf{c2} & \textbf{c3} &\textbf{c4}&\textbf{c5} \\ \hline  
   vin  5/4   & 1           & 1           & 2           & 0         & 1\\
   plat 5/3  & 1           & 0           & 0           & 1         & 2\\
   décor 5/2  & 0           & 1           & 0           & 1         & 0\\
   dessert 5/3  & 0           & 1           & 0           & 1         & 1\\\hline
   \end{tabular}
\end{center}
\end{solution}

  \item Donnez les normes des vecteurs représentant ces commentaires.
\begin{solution}
Les normes
\begin{itemize}
  \item $||c_1|| = \sqrt{1+ 1}$; $||c_2|| = \sqrt{3}$; $||c_3|| = \sqrt{4}$; 
    $||c_4|| = \sqrt{3}$, $||c_5|| = \sqrt{6}$  
\end{itemize}
\end{solution}
  
  \item Donner les résultats classés par similarité cosinus basée sur les tf 
  (on ignore l'idf) pour la requête associant les deux termes \texttt{décor} et 
  \texttt{dessert}.
     Expliquez brièvement la raison  du classement pour le premier document.

\begin{solution}
Les cosinus (requête non normalisée).

\begin{itemize}
  \item  c2 : $\frac{1+1}{\sqrt{3}}$
   \item  c4 :  $\frac{1+1}{\sqrt{3}}$  
    \item c5 : $\frac{1}{\sqrt{6}}$ 
 \end{itemize}

\end{solution}  

   \item Pour les documents c2 et c4,
   qu'est-ce qui change si on ne met pas de restriction 
   sur le vocabulaire (tous les mots sont indexés)?
   
\begin{solution}
Dans ce cas le plus petit document (c2) l'emporte car sa norme est
plus petite. Intuititvement: il parle plus \emph{spécifiquement} de
décor et de dessert que c4.
\end{solution}   
\end{enumerate}


\section{Systèmes distribués (3 pts)}

Rappelons que Elastic Search partitionne les documents avec une méthode 
de hachage statique, alors que 
Cassandra s'appuie sur la technique
du \textit{consistent hashing}.

On suppose que notre grappe de serveur comprend $n$ serveurs. 

\begin{enumerate}
  \item Expliquez comment un document de notre collection A de documents 
     JSON est inséré dans une grappe ElasticSearch (donnez l'algorithme 
     déterminant le serveur). 
     \item Même question pour une ligne de la table Cassandra \texttt{Notation}.
     \item Notre collection double de taille en quelques jours. Que peut-on 
     faire pour assurer la scalabilité avec chacune des deux solutions envisagées, 
     ElasticSearch et Cassandra\,?
\end{enumerate}

\end{document}
